% ============================================================
\chapter*{Pr\'eface}
\addcontentsline{toc}{chapter}{Pr\'eface}
% ============================================================

\section*{Pourquoi ce cours~?}

L'apprentissage profond a r\'evolutionn\'e de nombreux domaines de l'intelligence
artificielle~: vision par ordinateur, traitement du langage naturel, reconnaissance
vocale. Cependant, les architectures classiques --- r\'eseaux de neurones
enti\`erement connect\'es, r\'eseaux convolutifs sur grilles r\'eguli\`eres --- reposent
sur des hypoth\`eses structurelles fortes~: les donn\'ees vivent sur des grilles
euclidiennes r\'eguli\`eres.

Or, une quantit\'e croissante de probl\`emes fondamentaux en sciences et en
ing\'enierie fait intervenir des donn\'ees \`a structure \textbf{non euclidienne}~:
\begin{itemize}
    \item \textbf{Graphes}~: r\'eseaux sociaux, mol\'ecules, r\'eseaux de transport,
          graphes de connaissances.
    \item \textbf{Vari\'et\'es}~: surfaces 3D, espaces de configurations en robotique,
          donn\'ees sur la sph\`ere (m\'et\'eorologie, astrophysique).
    \item \textbf{Nuages de points}~: donn\'ees LiDAR, reconstruction 3D,
          g\'eom\'etrie mol\'eculaire.
    \item \textbf{Complexes simpliciaux}~: mod\'elisation de relations d'ordre
          sup\'erieur, topologie alg\'ebrique appliqu\'ee.
\end{itemize}

L'\textbf{apprentissage g\'eom\'etrique profond} (\textit{Geometric Deep Learning}, GDL) est
le cadre unificateur qui \'etend les principes de l'apprentissage profond aux
domaines non euclidiens, en s'appuyant sur les \textbf{sym\'etries} et les
\textbf{invariances} intrins\`eques des donn\'ees.

\section*{Le principe unificateur~: les sym\'etries}

Le fil conducteur de ce cours est le \textbf{principe d'\'equivariance}. L'id\'ee
fondamentale, h\'erit\'ee de la physique th\'eorique et du programme d'Erlangen de
Felix Klein, est que les structures math\'ematiques int\'eressantes sont caract\'eris\'ees
par leurs \textbf{groupes de sym\'etries}.

\begin{intuition}[title=\textbf{Le programme d'Erlangen pour le deep learning}]
De m\^eme que Klein a propos\'e de classifier les g\'eom\'etries par leurs groupes
de transformation, l'apprentissage g\'eom\'etrique profond propose de classifier
les architectures de r\'eseaux de neurones par les sym\'etries qu'elles respectent~:
\begin{center}
\begin{tabular}{lll}
\toprule
\textbf{Domaine} & \textbf{Sym\'etrie} & \textbf{Architecture} \\
\midrule
Grille r\'eguli\`ere & Translation & CNN \\
Ensemble & Permutation & Deep Sets \\
Graphe & Permutation des n\oe uds & GNN \\
Sph\`ere & Rotations $\mathrm{SO}(3)$ & Spherical CNN \\
Espace 3D & $\mathrm{SE}(3)$ & SE(3)-Transformer \\
\bottomrule
\end{tabular}
\end{center}
\end{intuition}

\section*{Public vis\'e}

Ce cours s'adresse aux \'etudiants de \textbf{Master 2} et de \textbf{doctorat} en
math\'ematiques appliqu\'ees, informatique, ou physique, ayant des pr\'erequis en~:
\begin{itemize}
    \item \textbf{Alg\`ebre lin\'eaire}~: espaces vectoriels, valeurs propres,
          d\'ecomposition spectrale.
    \item \textbf{Apprentissage profond}~: r\'eseaux de neurones, r\'etropropagation,
          optimisation par descente de gradient.
    \item \textbf{Probabilit\'es et statistiques}~: mesures de probabilit\'e,
          esp\'erance, estimation.
    \item \textbf{Programmation Python}~: ma\^itrise de PyTorch ou d'un framework
          \'equivalent.
\end{itemize}

Des notions de th\'eorie des groupes, de g\'eom\'etrie diff\'erentielle et de topologie
alg\'ebrique sont souhaitables mais seront introduites dans le cours.

\section*{Organisation du cours}

Le cours est structur\'e en onze chapitres, progressant des fondations
math\'ematiques vers les applications et la recherche de pointe~:

\begin{enumerate}
    \item \textbf{Motivations}~: pourquoi la g\'eom\'etrie en apprentissage profond.
    \item \textbf{Th\'eorie des groupes}~: sym\'etries, repr\'esentations, \'equivariance.
    \item \textbf{R\'eseaux de neurones sur graphes (GNN)}~: passage de messages,
          agr\'egation.
    \item \textbf{Variantes de GNN}~: GAT, GIN, GraphSAGE.
    \item \textbf{M\'ethodes spectrales}~: laplacien de graphe, convolutions spectrales.
    \item \textbf{Apprentissage sur les vari\'et\'es}~: g\'eom\'etrie diff\'erentielle,
          op\'erateurs intrins\`eques.
    \item \textbf{R\'eseaux \'equivariants et invariants}~: formalisme g\'en\'eral.
    \item \textbf{Nuages de points}~: PointNet et ses extensions.
    \item \textbf{Applications}~: chimie, physique, biologie, vision 3D.
    \item \textbf{Connexions avec la TDA}~: homologie persistante, repr\'esentations
          topologiques.
    \item \textbf{Directions de recherche}~: tendances actuelles et probl\`emes ouverts.
\end{enumerate}

\section*{Conventions et notations}

\begin{center}
\begin{tabular}{ll}
\toprule
\textbf{Notation} & \textbf{Signification} \\
\midrule
$G = (V, E)$ & Graphe avec sommets $V$ et ar\^etes $E$ \\
$\mathbf{A}$ & Matrice d'adjacence \\
$\mathbf{D}$ & Matrice des degr\'es \\
$\mathbf{L} = \mathbf{D} - \mathbf{A}$ & Laplacien combinatoire \\
$\tilde{\mathbf{L}} = \mathbf{D}^{-1/2}\mathbf{L}\mathbf{D}^{-1/2}$ & Laplacien normalis\'e \\
$\mathbf{X} \in \R^{n \times d}$ & Matrice des attributs des n\oe uds \\
$\mathbf{h}_v^{(\ell)}$ & Repr\'esentation du n\oe ud $v$ \`a la couche $\ell$ \\
$\mathcal{N}(v)$ & Voisinage du n\oe ud $v$ \\
$\rho: G \to \mathrm{GL}(V)$ & Repr\'esentation d'un groupe $G$ \\
$\mathcal{M}$ & Vari\'et\'e diff\'erentiable \\
$T_p\mathcal{M}$ & Espace tangent en $p$ \\
$\norm{\cdot}$ & Norme \\
$\inner{\cdot}{\cdot}$ & Produit scalaire \\
$\sigma$ & Fonction d'activation (ReLU, etc.) \\
$\bigoplus$ & Agr\'egation (somme, moyenne, max) \\
\bottomrule
\end{tabular}
\end{center}

\section*{Le cadre 5G de Bronstein et al.}

Le cadre unificateur de l'apprentissage g\'eom\'etrique profond repose sur cinq
domaines fondamentaux, les \og 5G \fg~:

\begin{definition}[Les 5G]
Les cinq domaines de l'apprentissage g\'eom\'etrique profond sont~:
\begin{enumerate}
    \item \textbf{Grids} (Grilles)~: donn\'ees sur grilles r\'eguli\`eres (images, s\'eries temporelles).
    \item \textbf{Groups} (Groupes)~: sym\'etries globales agissant sur les donn\'ees.
    \item \textbf{Graphs} (Graphes)~: donn\'ees relationnelles avec invariance par permutation.
    \item \textbf{Geodesics} (G\'eod\'esiques)~: donn\'ees sur des vari\'et\'es courbes.
    \item \textbf{Gauges} (Jauges)~: fibr\'es et connexions pour le transport parall\`ele.
\end{enumerate}
\end{definition}

\begin{remarque}[Unification]
Chacun de ces domaines peut \^etre vu comme un cas particulier d'un cadre g\'en\'eral
o\`u un \textbf{groupe de sym\'etries} $\mathfrak{G}$ agit sur un \textbf{domaine}
$\Omega$ et o\`u l'on cherche des fonctions $f: \mathcal{X}(\Omega) \to \mathcal{Y}$
qui sont \textbf{\'equivariantes} par rapport \`a l'action de $\mathfrak{G}$~:
\[
    f(\rho_{\mathcal{X}}(g) \cdot x) = \rho_{\mathcal{Y}}(g) \cdot f(x),
    \quad \forall g \in \mathfrak{G}.
\]
\end{remarque}

\section*{Ressources et r\'ef\'erences principales}

Ce cours s'appuie sur plusieurs ouvrages et articles fondateurs~:
\begin{itemize}
    \item M.~M.~Bronstein, J.~Bruna, T.~Cohen, P.~Veli\v{c}kovi\'c,
          \textit{Geometric Deep Learning: Grids, Groups, Graphs, Geodesics,
          and Gauges}, 2021.
    \item T.~N.~Kipf et M.~Welling, \textit{Semi-Supervised Classification with
          Graph Convolutional Networks}, ICLR 2017.
    \item P.~W.~Battaglia et al., \textit{Relational Inductive Biases, Deep
          Learning, and Graph Networks}, 2018.
    \item C.~R.~Qi et al., \textit{PointNet: Deep Learning on Point Sets for
          3D Classification and Segmentation}, CVPR 2017.
    \item N.~Thomas et al., \textit{Tensor Field Networks}, 2018.
    \item V.~G.~Satorras et al., \textit{E(n) Equivariant Graph Neural Networks},
          ICML 2021.
\end{itemize}

\section*{Logiciels utilis\'es}

Les travaux pratiques utilisent principalement~:
\begin{itemize}
    \item \textbf{PyTorch} (\texttt{torch}) --- framework d'apprentissage profond.
    \item \textbf{PyTorch Geometric} (\texttt{torch\_geometric}) --- biblioth\`eque pour
          l'apprentissage sur graphes.
    \item \textbf{NetworkX} --- manipulation et visualisation de graphes.
    \item \textbf{Open3D} --- traitement de nuages de points 3D.
    \item \textbf{GUDHI / Ripser} --- calculs d'homologie persistante.
\end{itemize}

\begin{codeblock}[title=\textbf{V\'erification de l'environnement}]
\begin{minted}{python}
import torch
import torch_geometric
import networkx as nx

print(f"PyTorch version: {torch.__version__}")
print(f"PyG version: {torch_geometric.__version__}")
print(f"CUDA disponible: {torch.cuda.is_available()}")

# Test rapide : creation d'un graphe PyG
from torch_geometric.data import Data
edge_index = torch.tensor([[0, 1, 1, 2],
                            [1, 0, 2, 1]], dtype=torch.long)
x = torch.randn(3, 16)  # 3 noeuds, 16 attributs
data = Data(x=x, edge_index=edge_index)
print(f"Graphe: {data.num_nodes} noeuds, {data.num_edges} aretes")
\end{minted}
\end{codeblock}

\section*{Remerciements}

Ce cours a b\'en\'efici\'e des contributions de nombreux chercheurs et \'etudiants.
Nous remercions particuli\`erement les auteurs du \textit{Geometric Deep Learning
Blueprint} pour avoir \'etabli le cadre conceptuel unificateur qui structure
l'ensemble de cet enseignement.

\vspace{1cm}
\begin{flushright}
\textit{[Author Name]} \\
\textit{Mars 2026}
\end{flushright}
